% history-ar.tex — Arabic edition of ch. 3, the history of sourdough
%
% Original: The Sourdough Framework by Hendrik Kleinwächter,
% the-sourdough-framework.com, CC BY-SA 4.0. The Arabic edition inherits
% CC BY-SA 4.0.
%
% Arabic layer of upstream history/sourdough-history.tex (scaffold
% structure, \label{ch:history}, figure labels, \cite keys and the
% quoting epigraph verbatim; text re-typeset from the recovery reference
% text/03-history.txt, Phase 28R unit 5 — no retranslation). Latin terms
% ride \textenglish{} islands with the parens inside, per the preface
% convention; \num{} keeps the two grouped years (12 000 / 60 000 BC) as
% the accepted edition's text layer shows them; \ch{CO2} rides a \LR{}
% island (the \pHvalue ch. 7 precedent: bare chemistry macros mirror
% their letters in RTL). The two TikZ timelines are the re-derived
% Arabic sources in figures-ar/ (R1-FIGURES route); einkorn-crumb and
% sourdough-stove are upstream artwork and carry over unchanged.
%
% Short captions: the recovery reference carries no short forms (they
% never render outside a list of figures, which no unit driver builds).
% The optional arguments below are minimal derivations from the full
% captions, flagged here and in status.md for the R3 LOF reconciliation.

\chapter{تاريخ الخبز الحامض}%
\label{ch:history}
\begin{quoting}
    نبدأ هذا الكتاب بالحديث باختصار عن تاريخ الخبز الحامض الطويل منذ القِدم، وكيف استعمل الناس عمليات مشابهة في أغذية أخرى كالجِعة. ثم اكتشاف الخميرة، وكيف أحدثت مع تطوّر الآلات ثورةً في صناعة الخبز. وفي الزمن الأقرب، تكوّنت مجتمعات حول الخبز الحامض والخَبز المنزلي تحاول إعادة تعلّم دروس الماضي.
\end{quoting}

تبدأ حكاية الخبز الحامض في محيطات ما قبل التاريخ، وهي المحيطات التي وُلدت فيها أشكال الحياة كلها على الأرض. ولتصوّر التاريخ العريض لكوكبنا، لنرسم خطاً زمنياً في سنة واحدة من 365 يوماً. على هذا المقياس، يدل الأول من يناير على تشكّل الأرض قبل 4.54 مليار سنة، ومنتصف ليل 31 ديسمبر هو الحاضر. كل يوم يمثّل نحو 12 مليون سنة. تبسّط هذه التقنية تعقيدات الزمن، لكنها تحوّل أيضاً المدى الخيالي من تاريخ كوكبنا إلى إطار زمني أمكن إدراكه. نحن البشر في الواقع إضافة حديثة إلى هذا الكوكب، حداثةً لدرجة أن ظهورنا الأول كان في مساء 31 ديسمبر. ويبدو أن البشر قد وصلوا في اللحظة المناسبة تماماً ليلحقوا بالاحتفال في نهاية السنة.

في 25 مارس ولّدت المحيطات أول البكتيريا وحيدة الخلية. وفي هذه المياه ازدهرت أيضاً كائنات حية أخرى وحيدة الخلية، هي العتائق \textenglish{(archaea)}، وهي كائنات تسكن البيئات القاسية، من الفتحات البركانية المغلية إلى المياه الجليدية.

\begin{figure}[!htb]
  \centering
  % fig-life-planet-sourdough-timeline-ar.tex — Arabic edition of the
% life/planet/sourdough timeline (ch. 3, figure 3.1)
%
% Original: The Sourdough Framework by Hendrik Kleinwächter,
% the-sourdough-framework.com, CC BY-SA 4.0. The Arabic edition inherits
% CC BY-SA 4.0.
%
% R1-FIGURES route: re-derived from the upstream post-#492 source
% figures/fig-life-planet-sourdough-timeline.tex (geometry, event
% positions, branch lines and the December cluster verbatim), with the
% label text recovered from the accepted edition's render of p. 15
% (book-ar-fixed.pdf, b78c2a28…) — label recovery against a visual
% reference, not retranslation. Arabic label text is wrapped in
% \textarabic{}: inside a tikzpicture the sty's env hook sets LTR, and
% bare Arabic node text renders un-reordered there (probe-verified,
% 2026-11-15); \textarabic restores the RTL base and the digit runs read
% LTR inside it exactly as in the accepted edition, including the
% parenthesized clock times. Latin and digit islands inside labels ride
% \LR{} rather than \textenglish{}: bare parens in the RTL \textarabic
% context mirror (probe-verified: (eukaryotes) and (23:59:59) rendered
% as )eukaryotes( and )23:59:59( before the \LR{} wrap).
% CORRECTION (probe5, 2026-11-15, real environment): the earlier claim
% that TikZ node interiors never pass the post_linebreak_filter is
% WRONG — arabic-digits.lua's remap_glyphs recurses into hlist/vlist
% heads unconditionally, so it DOES reach node interiors. Under the
% standing, owner-authorized remap (R2 amendment: every Amiri digit and
% ASCII letter the filter reaches is remapped to size-matched Pagella
% regardless of nesting), figure digits and figure Latin therefore ride
% TeXGyrePagella-Regular in this rebuild, where the accepted edition's
% SVG-route XObjects carried Amiri-only subsets. No content-layer fix
% exists: any Amiri the filter sees is remapped, and the lua is a
% no-touch-without-authorization file. Recorded as a route consequence
% for the owner: either accept the uniform Pagella digit/Latin face in
% figures (recommended; same policy as prose), or authorize a lua
% exemption amendment for figure boxes.

\begin{tikzpicture}
  % Draw horizontal line
  \draw[line width=1pt] (0,0) -- (\textwidth,0);

  % Define the width of each segment
  \pgfmathsetlengthmacro{\segmentwidth}{\textwidth/12}

  % Draw lines for the events, higher up so that they don't overflow the text
  % Placing the lines has been a bit manual work of trying different values
  % Maritime bacteria.

  \draw[line width=1pt] (2.8*\segmentwidth,1) -- (2.8*\segmentwidth,0.2);
  % Eukaryotes
  \draw[line width=1pt] (5.8*\segmentwidth,1.5) -- (5.8*\segmentwidth,0.2);
  % First bacteria on land
  \draw[line width=1pt] (9.1*\segmentwidth,-1.25) -- (9.1*\segmentwidth,-0.2);
  % Maritime fungi ancestors
  \draw[line width=1pt] (9.5*\segmentwidth,-2) -- (9.5*\segmentwidth,-0.2);
  % Fungi on land
  \draw[line width=1pt] (10.8*\segmentwidth,-2.75) -- (10.8*\segmentwidth,-0.2);
  % Yeasts on land
  \draw[line width=1pt] (11.1*\segmentwidth,-3.0) -- (11.1*\segmentwidth,-0.2);
  % First dinosaurs
  \draw[line width=1pt] (11.4*\segmentwidth,0.5) -- (11.4*\segmentwidth,0.2);
  % Pangea begins to rift apart
  \draw[line width=1pt] (11.6*\segmentwidth,1) -- (11.6*\segmentwidth,0.2);
  % Dinosaur extinction
  \draw[line width=1pt] (11.9*\segmentwidth,1.5) -- (11.9*\segmentwidth,0.2);
  % Additional line for dinosaurs since it is so close
  \draw[line width=1pt] (11.9*\segmentwidth,1.49) -- (11.70*\segmentwidth,1.85);

  % Special lines for december events since they are so close together
  \draw[line width=1pt] (12.0*\segmentwidth,3.0) -- (12.0*\segmentwidth,0.2);  % Main branch
  \draw[line width=1pt] (12.0*\segmentwidth,3.0) -- (11.75*\segmentwidth,2.5); % Branch to first humans
  \draw[line width=1pt] (12.0*\segmentwidth,3.0) -- (11.75*\segmentwidth,3.0); % Branch to Jordan
  % Move Pasteur down a bit so the lines look like they cross
  \draw[line width=1pt] (12.0*\segmentwidth,2.99) -- (11.75*\segmentwidth,3.5); % Branch to Pasteur

  % Draw months and month separators
  \foreach \i/\month in {0/{يناير}, 1/{فبراير}, 2/{مارس}, 3/{أبريل},
                         4/{مايو}, 5/{يونيو}, 6/{يوليو}, 7/{أغسطس},
                         8/{سبتمبر}, 9/{أكتوبر}, 10/{نوفمبر}, 11/{ديسمبر}} {
      % Separators
      \draw[line width=1pt] (\i*\segmentwidth,0.1) -- (\i*\segmentwidth,-0.1);
      % Month names
      \node[timeline_event, below] at ({(\i+0.5)*\segmentwidth},-0.1) {\textarabic{\month}};
  }
  \draw[line width=1pt] (\textwidth,0.1) -- (\textwidth,-0.1);

  % Full timeline width for billion years
  \draw[stealth-stealth, line width=1pt] (0,-3.8) -- node[midway, timeline_timespan] {\textarabic{4.45 مليار سنة}} (\textwidth,-3.8);

  % Indicator for the period of 3 months = 1.1 billion years
  \draw[stealth-stealth, line width=1pt] (0,-1.0) -- node[midway, timeline_timespan] {\textarabic{1.11 مليار سنة}} ({\segmentwidth * 3},-1.0);

  % Place events on the timeline with dates using the timeline_event style
  % As a calculation I used (4.54 billion years / 12 months = 0.3785 billion years/month.
  \node[timeline_event, above] at (2.5*\segmentwidth,1) {\textarabic{25 مارس: أول البكتيريا والعتائق البحرية}};
  \node[timeline_event, above] at (4.50*\segmentwidth,1.5) {\textarabic{25 يونيو: أول الكائنات ذات النواة \LR{(eukaryotes)}}};
  \node[timeline_event, above] at (7.8*\segmentwidth,-1.5) {\textarabic{4 أكتوبر: أول البكتيريا على اليابسة}};
  \node[timeline_event, above] at (8.0*\segmentwidth,-2.25) {\textarabic{15 أكتوبر: أول أسلاف الفطريات المائية}};
  \node[timeline_event, above] at (9.7*\segmentwidth,-2.75) {\textarabic{24 نوفمبر: الفطريات على اليابسة}};
  \node[timeline_event, above] at (10.5*\segmentwidth,-3.25) {\textarabic{3 ديسمبر: الخمائر على اليابسة}};
  \node[timeline_event, above] at (10.2*\segmentwidth,0.5) {\textarabic{14 ديسمبر: أول الديناصورات}};
  \node[timeline_event, above] at (9.8*\segmentwidth,1) {\textarabic{17 ديسمبر: بدء تفكك بانجيا}};
  \node[timeline_event, above] at (10.15*\segmentwidth,1.5) {\textarabic{29 ديسمبر: انقراض الديناصورات}};
  \node[timeline_event, above, anchor=east, align=right] at (11.75*\segmentwidth,2.5) {\textarabic{31 ديسمبر: أول البشر}};
  \node[timeline_event, above, anchor=east, align=right] at (11.75*\segmentwidth,3.0) {\textarabic{31 ديسمبر: أصل الخبز الحامض في الأردن \LR{(23:59:55)}}};
  \node[timeline_event, above, anchor=east, align=right] at (11.75*\segmentwidth,3.5) {\textarabic{31 ديسمبر: عزل لوِي باستور للخميرة \LR{(23:59:59)}}};

\end{tikzpicture}

  \caption[خط زمني للأحداث المهمة]{خط زمني للأحداث المهمة بدءاً من اليوم الأول لوجود الأرض، مقسّماً إلى أشهر، وممتداً حتى يومنا هذا المعلَّم عند منتصف الليل. تصوّر هذه الرسمة الخطوات المحورية في تاريخ الحياة والخبز الحامض على الأرض.}%
  \label{fig:planet-timeline}
\end{figure}

ما زال الأمر محل جدل: أيّهما جاء أولاً، البكتيريا أم العتائق؟ وعلى مدى ثلاثة أشهر (نحو 1.1 مليار سنة) هيمنت هذه الكائنات على المحيطات. ثم في 25 يونيو، وفي حدث بالغ الاستبعاد، ابتلعت عتيقةٌ بكتيرياً؛ وبدلاً من هضمها تكوّنت بينهما علاقة تعايش، فنشأت أول الكائنات ذات النواة، في محطة تطورية فارقة. من هذا الحدث تفرّعت النباتات والفطريات، وفي نهاية المطاف البشر أيضاً.

ظلّت الحياة مائية ثلاثة أشهر أخرى. في 4 أكتوبر استوطنت البكتيريا اليابسة لأول مرة، وحلول 15 أكتوبر ظهرت أول الفطريات المائية، ثم تأقلمت واستوطنت اليابسة بحلول 24 نوفمبر.

وبحلول 3 ديسمبر ظهرت الخمائر على اليابسة، فوُضعت بذلك لمسات الأساس لصناعة الخبز. ولنتخطَّ 140 مليون سنة إلى 14 ديسمبر: هنا نشأت الديناصورات. وبعد يومين فقط من ظهورها، في 17 ديسمبر، بدأت القارة العظمى بانجيا (\textenglish{Pangea}) بالتفكك، فأعادت تشكيل القارات إلى صورتها الحالية. بقيت الديناصورات في العرش حتى 29 ديسمبر حين أصابها الانقراض. وبعد 25 مليون سنة أخرى، أي بعد يومين من انقراضها في خط زمننا، ظهر البشر.

قبل ساعات قليلة من وصول البشر، كانت ثورةٌ طهوية أدقّ تجري على مهل. وبحلول \num{12000} قبل الميلاد، قبل 5 ثوانٍ فقط من منتصف ليلنا المجازي، كانت أول أرغفة الخبز الحامض تُخبز في الأردن القديم. وفي طرفة عين، أو 4 ثوانٍ في ضغطنا الزمني، مهّد عمل باستور الرائد مع الخمائر الطريق لصناعة الخبز الحديثة. ومن اللحظة التي بدأ فيها إعداد هذا الكتاب إلى لحظة قراءتك له الآن، لم يمرّ سوى أجزاء من الثانية~\cite{Yong+2017}.

وإذا تعمّقنا في عالم الخبز الحامض، فيُرجَّح أن يعود أثره إلى الأردن القديم المذكور آنفاً~\cite{jordan+bread}. وعلى مقياس تاريخ الأرض، يمكن اعتبار الخبز الحامض اختراعاً حديثاً جداً.

\begin{figure}[!htb]
  \centering
  % fig-sourdough-history-timeline-ar.tex — Arabic edition of the sourdough
% history timeline (ch. 3, figure 3.2)
%
% Original: The Sourdough Framework by Hendrik Kleinwächter,
% the-sourdough-framework.com, CC BY-SA 4.0. The Arabic edition inherits
% CC BY-SA 4.0.
%
% R1-FIGURES route: re-derived from the upstream post-#492 source
% figures/fig-sourdough-history-timeline.tex (axis geometry, the
% millenary/century tick loops with \num, the axis break and the era
% spans verbatim), with the label text recovered from the accepted
% edition's render of p. 16 (book-ar-fixed.pdf, b78c2a28…) — label
% recovery against a visual reference, not retranslation. Arabic label
% text is wrapped in \textarabic{} (see the companion figure 3.1 file for
% the probe evidence); the axis tick numbers stay bare \num{} — digits in
% the LTR picture base read correctly, as the accepted edition shows.
% The upstream Jordan label's two-line tabular is set as an align=right
% node instead: RTL lines flush right reproduce the accepted edition's
% render (second line under the right edge of the first), which an LTR
% l-column tabular cannot.
%
% Figure digit/Latin face (probe5, 2026-11-15): the standing remap in
% arabic-digits.lua reaches TikZ node interiors (it recurses into
% hlist/vlist heads unconditionally), so digits and Latin inside
% \textarabic{}/\LR{} islands ride TeXGyrePagella-Regular here, where
% the accepted edition's SVG-route XObjects carried Amiri-only subsets.
% Route consequence, documented in the companion figure 3.1 file; owner
% ruling pending (accept uniform Pagella, or authorize a lua exemption).

\begin{tikzpicture}
  \pgfmathsetlengthmacro{\timelinewidth}{(\textwidth-0.5cm)}

  % Define the width of each segment
  \pgfmathsetlengthmacro{\segmentwidth}{\timelinewidth/12}

  % Draw horizontal lines
  \draw[line width=1pt, color=hlorange] (0,0) -- (\timelinewidth/2,0);
  \draw[line width=1pt] (\timelinewidth/2,0) -- (\timelinewidth/2 +1.5*\segmentwidth/3,0);
  \draw[line width=1pt] (\timelinewidth/2 + 2*\segmentwidth/3,0) --(\timelinewidth, 0);
  \draw[line width=1pt] (\timelinewidth/2+1.5*\segmentwidth/3-5,-0.2) -- (\timelinewidth/2+1.5*\segmentwidth/3+5, 0.2);
  \draw[line width=1pt] (\timelinewidth/2+2*\segmentwidth/3-5,-0.2) -- (\timelinewidth/2+2*\segmentwidth/3+5, 0.2);

  % Lines for periods
  \draw[stealth-stealth, line width=1pt] (0,-3.7)
    -- node[midway, timeline_timespan] {\textarabic{صناعة الخبز التاريخية}} ({\segmentwidth * 7.8},-3.7);
  \draw[stealth-stealth, line width=1pt] ({\segmentwidth * 7.8},-3.7)
    -- node[midway, timeline_timespan] {\textarabic{الخبز الحديث}} ({\segmentwidth * 12},-3.7);

  % Regularly placed events, not in chronological order
  % since should be placed on top of others on the timeline

  % BC
  \draw[line width=1pt] ({\segmentwidth*3},1.0) -- ({\segmentwidth*3},0.3)
    node[at start, left, timeline_event] {\textarabic{6000 قبل الميلاد: أول جِعة في مصر}};
  \draw[line width=1pt] ({\segmentwidth*5.95},1.5) -- ({\segmentwidth*5.95},0.3)
    node[at start, left, timeline_event] {\textarabic{70 قبل الميلاد: أول طاحونة مائية}};

  % Sourdough in Jordan
  \draw[line width=1pt] (0,-0.3) -- (0,-1.5);
  \draw[line width=1pt] (0,-1.5) -- (0.25,-1.5);
  \node[timeline_event, below, anchor=west, align=right] at (0.25,-1.5)
    {\textarabic{12000 قبل الميلاد: الخبز الحامض في الأردن}\\
     \textarabic{زراعة القمح}};

  % AD
  \draw[line width=1pt] ({\segmentwidth*10.50},1.0) -- ({\segmentwidth*10.50},0.3)
    node[at start, above, timeline_event] {\hspace{1.0cm}\textarabic{1950: القمح الحديث}};
  \draw[line width=1pt] ({\segmentwidth*9.60},1.5) -- ({\segmentwidth*9.60},0.3)
    node[at start, above, timeline_event] {\textarabic{1868: الخميرة التجارية}};
  \draw[line width=1pt] ({\segmentwidth*7.8},2) -- ({\segmentwidth*7.8},0.3)
    node[at start, above, timeline_event] {\textarabic{1680: اكتشاف الكائنات الدقيقة}};

  \draw[line width=1pt] ({\segmentwidth*8.80},-1.25) -- ({\segmentwidth*8.80},-0.3)
    node[at start, left, timeline_event] {\textarabic{1785: الطاحونة البخارية}};
  \draw[line width=1pt] ({\segmentwidth*9.57},-1.75) -- ({\segmentwidth*9.57},-0.3)
    node[at start, left, timeline_event] {\textarabic{1857: عزل الخميرة}};
  \draw[line width=1pt] ({\segmentwidth*9.80},-2.25) -- ({\segmentwidth*9.80},-0.3)
    node[at start, left, timeline_event] {\textarabic{1885: الخلاط الكهربائي}};
  \draw[line width=1pt] ({\segmentwidth*11.20},-2.75) -- ({\segmentwidth*11.20},-0.3)
    node[at start, left, timeline_event] {\textarabic{2020: جائحة كورونا}};

  % Indicators for period
  % Draw millenary and century separators
  \foreach \i/\century in {0/-12000, 1/-10000, 2/-8000, 3/-6000, 4/-4000, 5/-2000}{
      % Separators
      \draw[line width=1pt, color=hlorange] (\i*\segmentwidth,0.1) -- (\i*\segmentwidth,-0.1);
      % Events for timeline
      \node[timeline_event, below, text=hlorange] at ({(\i)*\segmentwidth},-0.1) {\num{\century}};
  }
  \foreach \i/\century in {6/0, 7/1600, 8/1700, 9/1800, 10/1900, 11/2000, 12/2100}{
      % Separators
      \draw[line width=1pt] (\i*\segmentwidth,0.1) -- (\i*\segmentwidth,-0.1);
      % Events for timeline
      \node[timeline_event, below] at ({(\i)*\segmentwidth},-0.1) {\num{\century}};
  }
\end{tikzpicture}

  \caption[خط زمني لأهم الاكتشافات والأحداث]{خط زمني لأهم الاكتشافات والأحداث التي أدّت إلى الخبز الحامض الحديث.}%
  \label{fig:sourdough-timeline}
\end{figure}

تبقى الأصول الدقيقة للخبز المخمّر مجهولة؛ فمن أقدم أرغفة الخبز الحامض المحفوظة ما نُبش في سويسرا~\cite{switzerland+bread}.

\begin{figure}[ht]
  \centering
  \includegraphics[width=\textwidth]{einkorn-crumb}
  \caption[رغيف قرص العين القديم]{رغيف قرص العين (\textenglish{Einkorn}) القديم. لاحظ بنية اللُّب الكثيفة.}%
  \label{fig:einkorn-crumb}
\end{figure}

والقصة الشعبية الأخرى تحكي امرأةً في مصر تعجن عجين خبز قرب نهر النيل. نسيت المرأة العجين، وعند عودتها بعد أيام وجدته قد زاد حجمه وانبعث منها رائحة غريبة. قررت أن تخبزها على أي حال، فكان جزاؤها خبزاً أخفّ وأطرى وأطيباً طعماً بكثير. ومنذ ذلك اليوم واصلت صناعة الخبز على هذه الطريقة~\cite{egyptian+bread}.

لم يعلم الناس حينها أن كائنات دقيقة عجيبة الصغر هي سرّ تحسّن الخبز. ولا نعرف متى بدأوا بأخذ قليل من عجين اليوم السابق لعجين اليوم التالي. لكن بفعل ذلك وُلدت صناعة الخبز الحامض كما نعرفها اليوم: خمائر برية في الدقيق وفي الهواء، وبكتيريا تبدأ بتفكيك خليط الدقيق والماء. الخميرة تجعل العجين منتفخاً، والبكتيريا تمنحها حموضتها في المقام الأول. وهذه الكائنات الدقيقة المتمايزة تعمل في علاقة تعايش. وقد قدّر البشر البنية الهوائية المحسّنة والحموضة الخفيفة في العجين، كما امتدت مدة صلاحية هذا الخبز بفضل ارتفاع الحموضة.

وسرعان ما اكتُشفت عمليات مشابهة في تخمير الجِعة أو صناعة النبيذ؛ فكان يُؤخذ مقدار يسير من الدفعة السابقة ليُستعمل في الدفعة التالية. وهكذا خلق البشر خمائر الخبز الحديثة، وخمائر النبيذ، وخمائر الجِعة~\cite{egypt+beer}.

ومع كل دفعة على مر الزمن، أصبحت الخمائر والبكتيريا أقدر على هضم ما يُلقى إليها. فحين تُغذّي بادئتك أنت تنتقي بالتربية كائنات دقيقة تجيدي أكل دقيقك، ومع كل جيل يتعلم خبزُك الحامض تخميرَ الدقيق المتاح له بشكل أفضل. وهذا هو السبب\footnote{أمر مذهل إن فكرت فيه: طبّق الناس هذه العملية آلاف السنين دون أن يدروا ما يجري!} وراء ميل الخبز الحامض الأكثر نضجاً أحياناً إلى تخمير العجائن أسرع~\cite{review+of+sourdough+starters}. الخبز الحامض في ذاته علاقة تعايش، لكنه أيضاً تأقلم مع البشر وكوّن معهم علاقة تعايش: نحن نوفر له الغذاء والماء، فيجازينا بخبز لذيذ، وفي المقابل نؤويه ونحميه. تنتشر أبواغ البادئة عبر التلوث الهوائي أو الحشرات كذباب الفاكهة، وهذا يتيح للبادئة نشر أبواغها أبعد وأبعد حول العالم.

تشير الأدلة إلى طحن الحبوب مبكراً في شمالي أستراليا نحو \num{60000} قبل الميلاد، ولا سيما في المأوى الصخري مادجيديبي \textenglish{(Madjedbebe)} في أرض أرنهيم \textenglish{(Arnhem Land)}~\cite{aboriginal+grinding+stones}. لكن التقدم الأهم جاء لاحقاً، كما وثّقه الجغرافي الإغريقي القديم سترابون سنة 71 قبل الميلاد. وصف سترابون في كتاباته أول طاحونة حجرية مائية، تُعرف بالطاحونة الحبيبية \textenglish{(gristmill)}. رفعت هذه الطواحين إنتاج الدقيق من بضعة كيلوغرامات إلى عدة أطنان مترية في اليوم~\cite{history+mills}.

استعملت هذه الطواحين المبكرة عجلات مجاديف أفقية، سُميت في النهاية العجلات النوردية \textenglish{(Norse wheels)} لكثرة انتشارها في إسكندنافيا. تتصل العجلات بمحور، ويتصل المحور بدوره بحجر الطحن المتحرك المركزي. تدفع المياه الجارية العجلات، فتنتقل قوة الطحن إلى الحجر الثابت \textenglish{(bed)}، وهو حجر في حجمٍ وشكلٍ شبيهين بحجر الطحن. تصميم بسيط لا يحتاج إلى تروس، لكن له قيداً: سرعة دوران الحجر تعتمد على كمية المياه وسرعة جريانها، وهو ما يجعلها أنسب ما تكون لمناطق الجداول سريعة الجريان، التي تكثر في المناطق الجبلية~\cite{mills+scandinavia}.

في سنة 1680 قدّم عالمٌ استثنائي اسمه أنطوني فان ليفينهوك ابتكاراً رائداً غيّر إلى الأبد فهمنا للعالم المجهري، وفي نهاية المطاف لصناعة الخبز. كان فان ليفينهوك سيّد صناعة العدسات، خُلِق بسعي لا يشبع نحو العوالم التي تخفى على العين المجردة، ومن عمله الريادي وُلد أول مجهر حديث. ما ميّزه هو الجودة الاستثنائية لعدساته، القادرة على تكبير الكائنات الدقيقة العجيبة الصغر بمقدار 270 مرة. وحملته فضوليته التي لا تهدأ إلى رحلة استطلاع: فحص الذباب، ودرّس الشعر المسوّس، ثم أخيراً صوّب نظره نحو المياه الهادئة لبحيرة صغيرة قرب دلفت.

في هذا الموطن المائي الهادئ صنف مشاهدات مدهشة: اكتشف طحالب وكائنات راقصة دقيقة الصغر كانت خافية حتى ذلك الحين عن إدراك البشر. ولما أراد أن يشارك مجتمعه العلمي مكتشفاته المبشرة، لقي الشك، إذ كان صعب التصديق أن أحداً قد رأى آلاف الكائنات الراقصة الدقيقة، كائنات صغيرة لدرجة أن العين البشرية لا تدركها.

ولم يُثنِه الشك، فتابع مطاردته لما خفي، وصوّب عدسته نحو رواسب الجِعة المتروكة عند صانع جِعة. وهناك، في ذلك الوسط المعتم، صنع فان ليفينهوك التاريخ بوصفه أول إنسان ترى عيناه بكتيرياً وخميراً، فكشف عالماً كان مستتراً من قبل، وسيعيد تشكيل فهمنا لعلم الأحياء الدقيقة~\cite{Yong+2017+Leeuwen}.

وفي الوقت نفسه بدأ صانعو الجِعة يجرّبون استعمال البقايا الطينية الناتجة عن تخمير الجِعة لبدء صناعة العجائن. لاحظوا أن عجائن الخبز الناتجة أصبحت منتفخة، وأنها تفتقر مقارنةً بعملية الخبز الحامض إلى الحموضة في المنتج النهائي. ومن الأمثلة المشهورة ما ورد في تقرير سنة 1875: كتب إبن نورتن هورسفورد عن الخبز الصغير الشهير كايزر سِمِلن \textenglish{(Kaiser Semmeln)}، أي خبز الإمبراطور الصغير. هذه في جوهرها أرغفة صغيرة تُخمَّر بخميرة الجِعة بدلاً من البادئة الحامضة. ولأن العملية أغلى كلفة، فقد انتهى أمر أرغفة كهذه إلى أن يستهلكها رجال الطبقة النبيلة في فيينا~\cite{vienna+breadrolls}.

وقد مكّن تصنيع عملية الطحن الحبيبي، المبدوء في أواخر القرن الثامن عشر مع أوليفر إيفانز (1785) وتصاميمه لطواحين تنتج الدقيق بصفة متصلة دون يداً عاملة~\cite{evans+mill}، وتطوّرها لاحقاً إلى طواحين بخارية، من تحقيق تقدم كبير في إنتاج الخبز.

\begin{figure}[ht]
  \centering
  \includegraphics[width=\textwidth]{sourdough-stove}
  \caption{رغيف يُخبز فوق الموقد دون فرن.}%
  \label{fig:sourdough-stove}
\end{figure}

وحدث أكبر قفزة في صناعة الخبز الصناعي سنة 1857 حين اكتشف عالم الأحياء الدقيقة الفرنسي لويس باستور عملية التخمير الكحولي. أثبت أن الكائنات الدقيقة الخميرة هي سبب التخمير الكحولي لا محفزات كيميائية أخرى. وواصل بحثه وكان أول شخص يعزل سلالات خميرة نقية وينمّيها. وبعدها بقليل، في 1868، كان الأخوان شارل ومكسيميليان فلايشمان أول من سجّل براءة اختراع لسلالات خميرة نقية لصناعة الخبز، وقد عُزلت الخمائر المعروضة من دفعات خبز حامض. وبحلول 1879 بُنيت الآلات لتكثير الخميرة في طوارد مركزية كبيرة~\cite{fleischmann+history}. ثبت أن الخميرة النقية ممتازة في تخمير العجائن وبسرعة فائقة: ما كان يستغرق 10 ساعات لتخمير عجين خبز أصبح ممكناً في ساعة واحدة، فصارت العملية أكثر كفاءة بكثير. وما جعل في النهاية صناعة دفعات عجين ضخمة ممكنة هو اختراع العجّان الكهربائي. كثيراً ما يُنسب إلى روفوس إيستمان، المخترع الأمريكي، تقدم مهم في تكنولوجيا الخلاطات: ففي 1885 حصل على براءة اختراع لخلاط كهربائي مزوّد بآلية يدوية تُدار بمقبض مرفقي. لم يكن هذا الجهاز متقدماً أو منتشراً بقدر الخلاطات الكهربائية اللاحقة، لكنه كان محاولة مبكرة لمكننة عمليتي الخلط والعجن في المطبخ بالكهرباء. مثّل اختراع إيستمان خطوة مهمة في تطور الخلاطات الكهربائية، لكنه لم يبلغ رقيّ طرازات لاحقة كخلاط كيتشين إيد \textenglish{(KitchenAid)}، الذي طُرح سنة 1919 ويُعدّ غالباً من أول الخلاطات الكهربائية الناجحة انتشاراً، ولعب دوراً مهماً في إحداث ثورة في أجهزة المطبخ للطهاة المنزليين~\cite{first+mixer}~\cite{kitchenaid+history}.

وفي الحرب العالمية الثانية طُوّرت أول خميرة جافة معلبة، وقد أتاحت في نهاية المطاف للمخابز وللخبازين المنزليين صنع الخبز أسرع وبثبات أكبر. وبفضل الخميرة النقية صار بناء آلات صناعة الخبز الصناعية ممكناً. فإذا حافظت على الحرارة نفسها والدقيق نفسه وسلالات الخميرة نفسها، صار التخمير قابلاً للاستنساخ بدقة. وهذا ما قاد في نهاية المطاف إلى نشوء المخابز العملاقة ومراكز مزج الدقيق. فالمخابز طلبت دقيقاً مطابقاً من سنة إلى سنة ليُخبز في آلاتها؛ ولهذا السبب لا يوجد أي دقيق تشتريه اليوم من السوبرماركت من مصدر واحد: إنه دائماً مخلوط لتحقيق المنتج نفسه تماماً عبر السنين.

أما القمح الحديث، وتحديداً الأصناف عالية الغلة المقاومة للأمراض الشائع زراعتها اليوم، فبدأ تطويره في منتصف القرن العشرين، وتُعرف هذه الفترة غالباً بالثورة الخضراء \textenglish{(Green Revolution)}.

وأحد الشخصيات المحورية في هذا التطوير هو العالم الأمريكي نورمان بورلوغ، الذي يُنسب إليه الفضل في تربية أصناف قمح عالية الغلة، وخاصة أصناف القمح القزم، المقاومة للأمراض والقادرة على الازدهار في بيئات متنوعة. عمله، الذي بدأ في أربعينيات القرن العشرين وواصل حتى ستينياته، لعب دوراً حاسماً في رفع إنتاج القمح حول العالم وتخفيف النقص الغذائي~\cite{green+revolution}.

ومع تسارع أزمنة التخمير تدهور طعم الخبز النهائي. فعميلة الإنبات التي تستحثها إنزيمات معينة أساسية لتكوين قوام أخفّ وقشرة أطيب طعماً، ولا يمكن تسريعها إلى ما لا نهاية. وسرعان ما بدأت المخابز بإضافة إنزيمات خارجية لبلوغ خصائص تشبه الخبز الحامض في العجائن المعتمدة على الخميرة. وكاد الخبز الحامض يختفي تماماً من وجه الأرض، فلم يبقَ على صنعه سوى قلة من المتحمسين الشغوفين.

وبدأ الناس فجأة يتحدثون أكثر عن داء السيلياك \textenglish{(coeliac disease)} ودور الغلوتين. المرض ليس جديداً؛ فقد وُصف لأول مرة سنة 250 ميلادية~\cite{coeliac+disease}. لاحظ الناس أن الخبز الحديث يحوي غلوتيناً أكثر بكثير من الخبز القديم، وأن الخبز قديماً كان على الأرجح أكثر تسطحاً بكثير. وقد رُبّيت الحبوب عبر الزمن أكثر فأكثر نحو احتوائها على كمية أعلى من الغلوتين. الغلوتين بروتين يمنح الخبز الحديث بنية اللُّب الناعمة المنتفخة المميزة له. ترتبط بروتينات الغلوتين ببعضها متى نشّطها الماء. وعلى مدى التخمير يُحتجز \LR{\ch{CO2}} في هذه المصفوفة البروتينية، فتتوسع الجيوب الدقيقة المتكوّنة أثناء الخَبز، وبما أن العجين يتحول إلى هلام عند تسخينه، تتصلّب البنية. هذا ما يجعل الخبز يبدو ناعماً منتفخاً في المذاق. وكما هو حال شرب حليب البقر الخام، قد يتفاعل جهازك المناعي مع البروتينات المتناولة: هناك حساسية الغلوتين، وهناك داء السيلياك. وحين يقول الناس إنهم لا يتقبلون الغلوتين جيداً فغالباً إنما يصفون حساسية الغلوتين. ويصف بعض الناس مشكلات مشابهة عند تناول كمية زائدة من اللاكتوز، أما إذا أكلت جبناً مخمّراً طويلاً فإن معظم اللاكتوز تكون قد خمّرته الكائنات الدقيقة العجيبة الصغر. وقد درس الناس ولاحظوا أن الخبز الحامض يُتقبّل عادةً بشكل أفضل من الخبز الصناعي السريع العادي. وسبب ذلك أن الإنزيمات تحتاج إلى وقت لتعمل على العجين. الغلوتين بروتين تخزيني في الدقيق، وبمجرد أن يحفّز الماءُ الإنبات، يبدأ إنزيم البروتييز بتحويل الغلوتين إلى أحماض أمينية أصغر يحتاجها الإنبات. ومع الوقت أنت عملياً تخسر الغلوتين لأنه يتفتت طبيعياً. وعلاوة على ذلك، كانت بكتيريا حمض اللبنيك تبدأ تقليدياً بتفكيك خليط الدقيق والماء؛ ففي الطبيعة يُعاد تدوير كل شيء تقريباً، وجزء من غذائها استهلاك البروتينات في العجين. الخبز الحديث أسرع ولم تعد تحوي بكتيريا حمض اللبنيك. والعاملان معاً يعنيان أنك تستهلك منتجات بقيمة غلوتين أعلى بكثير مقارنة بالزمن القديم حين كان التخمير طبيعياً~\cite{raffaella+di+cagno}.

وفي زمن حمى الذهب في كاليفورنيا حمل الخبازون الفرنسيون ثقافة الخبز الحامض إلى أمريكا الشمالية، فذاع خبزٌ شعبي هو الخبز الحامض السان فرانسيسكوني، ويتميز بطعمه الحامض الفريد (الذي كان في السابق شائعاً في كل خبز). لكنه بقي طعاماً ضيقاً الانتشار بينما كان الخبز الصناعي في صعود. وما سرّع فعلاً عودة الخبز الحامض هو جائحة كورونا \textenglish{(COVID-19)} في 2020. أصبح الدقيق والخميرة نادرَين في السوبرماركت. وعندما عاد الدقيق تعذّر العثور على الخميرة، فبدأ الناس بالبحث عن بدائل وأعادوا اكتشاف الطريقة القديمة لصنع الخبز الحامض. وسرعان ما أدرك كثيرون أن صنع الخبز الحامض أكثر تعقيداً من الخبز الحديث المعتمد على الخميرة: عليك الاعتناء ببادئة خبز حامض وإبقاؤها في حالة مثالية لتخمير عجينك على نحو سليم. كما أنه مقارنةً بعجين معتمد على الخميرة لا يمكنك أن تدك العجين وترك التخمير مواصلاً. ويمكنك أن تفرط في تخمير عجينك، فتظهر النتيجة عجيناً لزجاً فوضوياً. قاد هذا التعقيد خبازين كثيرين إلى طلب المساعدة، وتكوّنت حول موضوع الخبز المنزلي مجتمعات مزدهرة كثيرة.

وفي مقابلة مع كارل دي سميدت (مالك مكتبة الخبز الحامض \textenglish{Sourdough Library}) قال شيئاً غيّر طريقة تفكيري في الخبز: «مستقبل الخبز الحديث يكمن في الماضي~\cite{interview+karl+de+smedt}.»
